\subsection{Comparison of enrichment score rankings for the Favalli CAIX selection}
\label{supp:enrichment-scores-comparison}

To compare the ranking schemes implemented in DELT-Hit, we re-analysed the Favalli selection against CAIX. As a
literature-supported positive-control motif, we focused on building block 1, code 474, which occurs in previously
reported CAIX-binding compounds from related DEL selections\cite{bassi_singlestrandeddnaencodedchemical_2020,favalli_stereoregiodefineddnaencoded_2021}. For the count-based and edgeR analyses, compounds were ordered by the corresponding aggregate score, and the occurrence of building block 1/code 474 was counted among the top 10, 25, 50, and 100 ranked compounds. Because the z-score method is computed at the individual-selection level, compounds were ranked separately for each replicate, the occurrence of code 474 was counted independently within each replicate-specific top list, and the bar plot reports the mean with the replicate standard deviation. This recovery analysis provides a high-level comparison of how effectively the different ranking strategies prioritize a known CAIX-associated building block. The analysis indicates that the edgeR-based methods are more sensitive than pure count-based ranking. The top-50 recovery comparison is summarized in Supplementary Figure~\ref{fig:favalli-ca9-ds-top50-code1}. Scripts to reproduce the analysis are described in \path{supporting_material/experiments/favalli/README.md}.

\begin{figure}[htbp]
\centering
\includegraphics[width=0.9\linewidth]{\suppfigroot/figures/favalli-ca9-top50-code1-counts.png}
\caption{Occurrence counts of code 1 building blocks among the top 50 ranked compounds for the Favalli CAIX selection analysis. For the z-score method, bars show the mean across replicate-specific top-50 selections and the error bars show the replicate standard deviation.}
\label{fig:favalli-ca9-ds-top50-code1}
\end{figure}
